\documentclass[10pt,a4paper]{amsart}
\usepackage[margin=19mm]{geometry}
\usepackage{amsmath,amssymb,mathtools}
\usepackage[scr=rsfs,cal=euler]{mathalfa}
\usepackage{xfrac}
\usepackage[unicode,hidelinks,bookmarks=false]{hyperref}
\newcommand{\ie}{i.e.\ }
\newcommand{\cf}{cf.\ }
\newcommand{\resp}{respectively\ }
\newcommand{\Subsection}{\subsection}
\providecommand{\nobreakdash}{\nobreak\hbox{-}}
\setlength{\emergencystretch}{2em}
\multlinegap0pt
\usepackage{bm}
\usepackage{bbm}
\usepackage{dsfont}
\usepackage{xspace}
\usepackage{cancel}
\usepackage{mathtools}
\usepackage{amsthm}
\makeatletter
\@ifundefined{theorem}{\newtheorem{theorem}{Theorem}}{}
\@ifundefined{lemma}{\newtheorem{lemma}[theorem]{\rm\bf Lemma}}{}
\makeatother
\title[Projective rigidity of circle patterns and polyhedral surfac]{Projective rigidity of circle patterns and polyhedral surfaces in hyperbolic ends}
\author{Jean-Marc Schlenker}
\date{Source: arXiv:2508.15339v1 (CC BY 4.0)}
\begin{document}
\maketitle

\noindent\emph{Source and licence.} Projective rigidity of circle patterns and polyhedral surfaces in hyperbolic ends. Version arXiv:2508.15339v1 (CC BY 4.0), distributed under the Creative Commons Attribution 4.0 International licence, as recorded in the arXiv licence metadata for this exact version and shown on the abstract page https://arxiv.org/abs/2508.15339v1. Full rights notice: licenses/arxiv-2508.15339-CC-BY-4.0.txt.\par\smallskip

\noindent\textit{Editorial scope.} Counted unit: the Lemma whose source label is \texttt{lm:pak} in the article (source0.tex lines 690--692, source label \texttt{lm:pak}), delivered together with the article's own complete original proof of it (source0.tex lines 694--715, one \texttt{proof} environment of 22 lines). No mathematics is written, repaired or re-derived: statement and proof are the article's own text line for line. Required context retained uncounted: none. Every definition, display and label of those lines is retained unchanged. Omitted from this package and not redistributed: 1--689, 693--693, 716--1126. Nothing inside a retained excerpt is omitted or altered: every retained body is delivered exactly as the article's lines are written, so no omission receipt of any kind accompanies this package. Citations: the retained excerpts carry no citation command at all, so no bibliography entry of the article is transcribed and no reference is redistributed. Reference adaptation: none was needed and none was made; every reference inside the delivered unit resolves inside it, so no unresolved reference or citation is produced. Printed slips kept literal: no printed word, symbol, hypothesis or number of the counted statement or of its proof was altered. Layout and class: the article's own class is \texttt{amsart} with packages including mathrsfs, amsmath, amssymb, enumerate, epsfig, fancyhdr, color, epstopdf, amsthm, latexsym, amscd, psfrag, graphicx, epsf; the wrapper is the lane's pinned \texttt{amsart} editorial wrapper at 10pt on A4, which restates the theorem-like environments the retained text needs and adds the article's own macro definitions that the retained text needs (none), with extra wrapper packages (bm, bbm, dsfont, xspace, cancel, mathtools). The delivered numbering is the unit's own. Assets: no retained span contains a figure environment, a graphics-inclusion command, a PDF-page inclusion, a table or a TikZ picture; no image file is copied into this package, the package declares no asset dependency and no image file is redistributed with it. Licence: version arXiv:2508.15339v1 of this article is distributed under the Creative Commons Attribution 4.0 International licence, as recorded in the arXiv licence metadata for this exact version and shown on the abstract page https://arxiv.org/abs/2508.15339v1. A full rights notice is delivered at licenses/arxiv-2508.15339-CC-BY-4.0.txt. Characters: every retained excerpt is pure ASCII, so the delivered engine, XeLaTeX with its default Latin Modern text font, reports no missing character.
\begin{lemma}\label{lm:pak}
  Let $\Sigma$ be a closed oriented surface of genus at least $2$. The only tight decoration of any quasi-simplicial triangulation of $\Sigma$ is the trivial one, that is, the decoration for which no edge is oriented.
\end{lemma}

\begin{proof}
  First remove from $\Sigma$ all triangles where no edge is oriented. This leaves a (possibly disconnected) triangulated surface with boundary, for which all boundary edges non-oriented.  Consider a connected component of this surface with boundary. It is a surface of genus $g$ with $b$ boundary components, with $\chi=2-2g-b<0$. Let $e_b$ be the number of boundary edges. 

  Note that in this surface, each triangle has at least one oriented edge (since triangles with no oriented edge has been removed). It follows that each triangle has at least one change of orientation, indeed:
  \begin{itemize}
  \item if all three edges are oriented, the orientation cannot be compatible at all corners, so there is at least one change of orientation,
  \item if exactly one edge is non-oriented, there is half a change of orientation at each of its endpoints,
  \item if exactly one edge is oriented, there is half a change of orientation at each of its endpoints.
  \end{itemize}

  We now call $c$ the total number of changes of orientations. Then $c\leq 2|V|-e_b$, since by the definition of a tight decoration the total inversion at each boundary vertex is at most $1$, and $e_b$ is also the number of boundary vertices. Moreover, $c\geq |F|$ since any triangle with at least one oriented edge has at a sum of changes of orientation of at least $1$. 

  $$ |V|-|E|+|F| = 2-2g-b~, $$
  while
  $$ 3|F| = 2|E|-e_b~, $$
  because each interior edge bounds $2$ faces, while each boundary edge bounds only one face. As a consequence,
  $$ 2|V|-(3|F|+e_b)+2|F| = 4-4g-2b~, $$
  so
  $$ 2|V|-e_b = |F| + (4-4g-2b) < |F|~. $$

  This is a contradiction since $c\geq |F|$ while $c\leq 2|V|-e_b$.  
\end{proof}

\end{document}
